\section{Zadanie 8}
Niech \(X\) będzie zmienną losową mówiącą, ile razy spśród \(n\) niezależnych rzutów
kostką wypadła 6-tka. Ogranicz \(P(X \geq n/4)\) za pomocą nierówności Markowa,
Czebyszewa i Czernowa. Podaj wartości numeryczne dla \(n = 1000\).

\subsection*{Rozwiązanie}
Mamy \(X \sim \text{Bin}(n, p)\) gdzie \(p = 1/6\). Stąd \(E[X] = n/6\) i \(\variance{X} = np(1-p) = 5n/36\).

\textbf{Nierówność Markowa:}
Dla \(X \geq 0\) mamy \(P(X \geq a) \leq \frac{E[X]}{a}\). Więc:
\[
	P(X \geq n/4) \leq \frac{n/6}{n/4} = \frac{4}{6} = \frac{2}{3}
\]

\textbf{Nierówność Czebyszewa:}
Dla dowolnej zmiennej losowej \(P(|X - E[X]| \geq t) \leq \frac{\variance{X}}{t^2}\).
Zauważmy, że \(X \geq n/4\) implikuje \(X - n/6 \geq n/4 - n/6 = n/12\), więc:
\[
	P(X \geq n/4) \leq P(|X - n/6| \geq n/12) \leq \frac{5n/36}{(n/12)^2} = \frac{5n/36}{n^2/144} = \frac{20}{n}
\]

\textbf{Nierówność Czernowa:}
Dla \(X \sim \text{Bin}(n,p)\) i \(\delta > 0\) mamy \(P(X \geq (1+\delta)np) \leq e^{-\frac{\delta^2 np}{2+\delta}}\).
Tu \((1+\delta)np = n/4\), więc \((1+\delta) \cdot n/6 = n/4\), czyli \(\delta = 1/2\). Stąd:
\[
	P(X \geq n/4) \leq e^{-\frac{(1/2)^2 \cdot n/6}{2 + 1/2}} = e^{-\frac{n/24}{5/2}} = e^{-\frac{n}{60}}
\]

\subsection*{Odpowiedź dla \(n=1000\)}
\begin{align*}
	\text{Markov:}    & \quad P(X \geq 250) \leq \frac{2}{3} \approx 0.6667                          \\
	\text{Czebyszew:} & \quad P(X \geq 250) \leq \frac{20}{1000} = 0.02                              \\
	\text{Czernow:}   & \quad P(X \geq 250) \leq e^{-1000/60} = e^{-50/3} \approx 5.7 \times 10^{-8}
\end{align*}
\[
	\boxed{\text{Markov: } 2/3, \quad \text{Czebyszew: } 0.02, \quad \text{Czernow: } e^{-50/3} \approx 5.7 \times 10^{-8}}
\]
